\documentclass[12pt,a4paper,oneside]{article}
\usepackage[brazil]{babel}
\usepackage[utf8]{inputenc}
\usepackage{olymp}
\usepackage{color}
\usepackage[dvipsnames]{xcolor}
\usepackage{listings}

\setcounter{problem}{5}

\begin{document}

\begin{problem}{Soma de frações}{frac}{2}

Faça um programa para ler duas frações e escrever o resultado de sua soma. 

Seu programa deve obrigatoriamente usar uma função com o seguinte protótipo:

\texttt{\textcolor{YellowGreen}{\textbackslash\textbackslash Recebe frações na/da e nb/db e retorna por referência nr/dr = na/da + nb/db}}\\
\texttt{\textcolor{Mulberry}{void} somafrac(\textcolor{Mulberry}{int} na, \textcolor{Mulberry}{int} da, \textcolor{Mulberry}{int} nb, \textcolor{Mulberry}{int} db, \textcolor{Mulberry}{int \&}nr, \textcolor{Mulberry}{int \&}dr)};

Note que os quatro primeiros parâmetros representam numeradores e denominadores das frações passadas como argumento, e os dois últimos são referências ao numerador e ao denominador da fração resultante da soma.

% \Notes
% A fração resultante deve estar \textbf{simplificada}!


\InputFile

A entrada contém duas frações, dadas no formato  $A/B$\; $C/D$, sendo $A, B, C, D$ números inteiros. Restrições: $0 \leq A,C \leq 1000$, $1 \leq B,D \leq 1000$.

\OutputFile

Seu programa deve gerar apenas uma linha na saída, no formato ``\texttt{A/B+C/D=E/F}'', sendo $A, B, C, D$ os valores lidos na entrada e $E$ e $F$ respectivamente o numerador e denominador da fração resultante da soma.


\Notes
A fração resultante deve estar \textbf{simplificada}!

\Example

\begin{example}
\exmp{
1/2 1/3
}{
1/2+1/3=5/6
}%
\end{example}


\begin{example}
\exmp{
1/10 1/10
}{
1/10+1/10=1/5
}%
\end{example}

\begin{example}
\exmp{
0/5 100/300
}{
0/5+100/300=1/3
}%
\end{example}

\begin{example}
\exmp{
5/4 6/8
}{
5/4+6/8=2/1
}%
\end{example}

%Comentários sobre o exemplo, se necessário.

\end{problem}

\end{document}
